% ===========================================================================
\section{Projective Varieties}\label{sec:proj}
\warmth{0}\quad\whisper{Add points at infinity and the geometry gets simpler: parallel lines finally meet.}

\begin{strand}
Projective space $\PP^n=\PP(K^{n+1})$ is the set of \emph{lines through the origin} in $K^{n+1}$,
with coordinates $[a_0:\cdots:a_n]$ defined up to scaling. It is covered by $n+1$ affine charts
$\{a_i\ne0\}\cong K^n$, and over $\R,\C$ it is \keyword{compact}. Geometers prefer $\PP^n$: parallel
lines in $K^2$ do not meet, but any two lines in $\PP^2$ do. The price is that we may only use
\keyword{homogeneous} polynomials, since only their vanishing is well defined on a line.
\end{strand}

\block{Homogeneous ideals, cones, closures}

\begin{definition}[projective variety]
For homogeneous $f_1,\dots,f_k$, $\ \V(f_1,\dots,f_k)=\{[a]\in\PP^n: f_i(a)=0\ \forall i\}$. An ideal is
\keyword{homogeneous} if generated by homogeneous polynomials. The \keyword{affine cone} $\hat X\subseteq
K^{n+1}$ is cut by the same ideal, and
\[
  \dim(X)=\dim(\hat X)-\fillin[0.8cm],\qquad \deg(X)=\deg(\hat X).
\]
\end{definition}

\recall{Why is the cone $\hat X$ singular at the origin once $\deg\ge2$? Every line of the cone passes through $0$.}

The \keyword{projective closure} $\bar Y\subseteq\PP^n$ of an affine $Y\subseteq K^n$ has ideal $\bar I$
obtained by \keyword{homogenizing}, and Proposition 2.18 computes it cleanly: take a degree-compatible
reduced Gröbner basis $\mathcal G$ of $I$ and homogenize each $g\in\mathcal G$.

\begin{theorem}[invariants survive the closure, Cor.\ 2.19]\label{cor:closure}
$\dim(\bar Y)=\dim(Y)$ and $\deg(\bar Y)=\deg(Y)$.
\end{theorem}

\begin{yourturn}
(Rational normal curve, Ex.\ 2.20.) Let $I=\gen{x_i-x_1^i:i=2,\dots,n}$, so $Y=\V(I)$ is the curve
$t\mapsto(t,t^2,\dots,t^n)$. Its projective closure $\bar Y\subset\PP^n$ has ideal generated by the
$2\times2$ minors of
\[
  \begin{pmatrix} x_0 & x_1 & x_2 & \cdots & x_{n-1}\\ x_1 & x_2 & x_3 & \cdots & x_n \end{pmatrix}.
\]
How many such minors are there? $\fillin[1.5cm]$. The curve $\bar Y$ has degree $\fillin[1cm]$ and is
the \keyword{rational normal curve}. \recall{The $\binom n2$ minors; degree $n$. The initial ideal is $\gen{x_1,\dots,x_{n-1}}^2$.}
\end{yourturn}

\block{Why projective: intersections always behave}

\begin{theorem}[projective intersection, Thm.\ 2.22]\label{thm:intersect}
Let $K$ be algebraically closed and $X,Y\subseteq\PP^n$ projective varieties with $d_1=\dim X$,
$d_2=\dim Y$. Then
\[
  \dim(X\cap Y)\;\ge\; \fillin[2.5cm].
\]
In particular, if $d_1+d_2\ge n$ then $X\cap Y$ is \fillin[2.5cm].
\end{theorem}

\begin{yourturn}
(Ex.\ 2.23, why the hypotheses matter.) Two quadric surfaces $X,Y\subseteq\PP^3$ ($n=3$, $d_1=d_2=2$).
\begin{itemize}[leftmargin=1.3em,itemsep=1pt]
  \item Theorem prediction: $\dim(X\cap Y)\ge 2+2-3=\fillin[0.8cm]$, and the intersection is nonempty.
  \item Over $\C$: their intersection is a union of four lines, $\dim=1$. \checkmark
  \item Over $\R$: the book's pair meets in only \emph{two points}, $\dim=0$. Which hypothesis of the
        theorem just failed? \fillin[3cm]
\end{itemize}
\end{yourturn}

\begin{strand}
Many models in science are projective varieties given by homogeneous constraints. Two showcases:
the \keyword{nilpotent $n\times n$ matrices} form an irreducible projective variety in $\PP^{n^2-1}$ of
$\dim n^2-n-1$ and $\deg n!$, a complete intersection cut by the non-leading coefficients of the
characteristic polynomial (for $n=2$: $\Ideal(X)=\gen{\operatorname{trace}(A),\det(A)}$); and the
\keyword{Kalman varieties} of matrices with an eigenvector in a fixed subspace, basic in control theory.
The lesson of the chapter: when constraints are homogeneous, regard the model as living in $\PP^n$.
\end{strand}

\loose{Exercise 17 (book): $X\subset\C^4$ cut by two random cubics. What does Bézout tell you about $\dim,\deg$, and what does Bertini tell you about smoothness and primeness?}
